\documentclass[UTF8,11pt]{ctexart}
\usepackage[a4paper,margin=22mm]{geometry}
\usepackage{amsmath,amssymb,bm,graphicx,adjustbox,fvextra,longtable,array}
\xeCJKDeclareCharClass{CJK}{"2160 -> "217F, "2460 -> "249B, "2200 -> "22FF}
\newcommand{\fitmath}[1]{\begin{adjustbox}{max width=\linewidth}$\displaystyle #1$\end{adjustbox}}
\setlength{\parindent}{0pt}\setlength{\parskip}{.65em}\renewcommand{\arraystretch}{1.3}
\sloppy\allowdisplaybreaks
\begin{document}
\section*{2013 年考研政治真题}
\subsection*{2013 年全国硕士研究生入学统一考试思想政治理论试题}

\subsection*{一、单项选择题}

1.有一幅对联，上联是“桔子洲，洲旁舟，舟行洲不行”；下联“天心阁，阁中鸽，鸽飞阁不飞。”这形象地说明了运动和静止是相互依存的，静止是：


\begin{center}\begin{tabular}{>{\raggedright\arraybackslash}p{\dimexpr(\linewidth-8\tabcolsep)/4\relax}>{\raggedright\arraybackslash}p{\dimexpr(\linewidth-8\tabcolsep)/4\relax}>{\raggedright\arraybackslash}p{\dimexpr(\linewidth-8\tabcolsep)/4\relax}>{\raggedright\arraybackslash}p{\dimexpr(\linewidth-8\tabcolsep)/4\relax}}
\textbf{A.} 运动的衡量尺度 & \textbf{B.} 运动的内在原因 & \textbf{C.} 运动的普遍状态 & \textbf{D.} 运动的存在方式\\[.35em]\end{tabular}\end{center}

2.一位机械工程专家讲过这样一件事：“文革”中，他在农场劳动，有一天领导要他去割羊草。他没养过羊，怎么认得羊草呢？但脑子一转办法就来了。他把羊都赶去，看羊吃什么就割什么。不到半天就割回了羊草。这位专家之所以这样做是因为他意识到，“羊吃草”与“割羊草”两者之间存在着：


\begin{center}\begin{tabular}{>{\raggedright\arraybackslash}p{\dimexpr(\linewidth-8\tabcolsep)/4\relax}>{\raggedright\arraybackslash}p{\dimexpr(\linewidth-8\tabcolsep)/4\relax}>{\raggedright\arraybackslash}p{\dimexpr(\linewidth-8\tabcolsep)/4\relax}>{\raggedright\arraybackslash}p{\dimexpr(\linewidth-8\tabcolsep)/4\relax}}
\textbf{A.} 主观联系 & \textbf{B.} 必然联系 & \textbf{C.} 因果联系 & \textbf{D.} 本质联系\\[.35em]\end{tabular}\end{center}

3.《资本论》中有这样的表述“对上衣来说，无论是裁缝自己穿还是他的顾客穿，都是一样的。”这主要是因为无论谁穿：


\begin{center}\begin{tabular}{>{\raggedright\arraybackslash}p{\dimexpr(\linewidth-4\tabcolsep)/2\relax}>{\raggedright\arraybackslash}p{\dimexpr(\linewidth-4\tabcolsep)/2\relax}}
\textbf{A.} 上衣都起到着使用价值的作用 & \textbf{B.} 上衣都起到着价值的作用\\[.35em]
\textbf{C.} 上衣都是抽象劳动的结果 & \textbf{D.} 上衣都是社会劳动的结果\\[.35em]\end{tabular}\end{center}

4.某资本家投资100 万元，每次投资所获得的利润是15 万元，假定其预付资本的有机构成是4:1，那么该资本家每次投资所实现的剩余价值率为：


\begin{center}\begin{tabular}{>{\raggedright\arraybackslash}p{\dimexpr(\linewidth-8\tabcolsep)/4\relax}>{\raggedright\arraybackslash}p{\dimexpr(\linewidth-8\tabcolsep)/4\relax}>{\raggedright\arraybackslash}p{\dimexpr(\linewidth-8\tabcolsep)/4\relax}>{\raggedright\arraybackslash}p{\dimexpr(\linewidth-8\tabcolsep)/4\relax}}
\textbf{A.} 50\% & \textbf{B.} 75\% & \textbf{C.} 100\% & \textbf{D.} 125\%\\[.35em]\end{tabular}\end{center}

5.当今世界是开放的世界，中国的发展离不开世界。实行对外开放是我国的一项基本国策。坚持这一国策的基本立足点是：


\begin{center}\begin{tabular}{>{\raggedright\arraybackslash}p{\dimexpr(\linewidth-8\tabcolsep)/4\relax}>{\raggedright\arraybackslash}p{\dimexpr(\linewidth-8\tabcolsep)/4\relax}>{\raggedright\arraybackslash}p{\dimexpr(\linewidth-8\tabcolsep)/4\relax}>{\raggedright\arraybackslash}p{\dimexpr(\linewidth-8\tabcolsep)/4\relax}}
\textbf{A.} 内外联动，互惠互利 & \textbf{B.} 多元平衡，共同发展 & \textbf{C.} 相互借鉴，求同存异 & \textbf{D.} 独立自主，自力更生\\[.35em]\end{tabular}\end{center}

6.公益性文化事业是保障公民基本文化权益的重要途径。大力发展公益文化事业，始终坚持放到首位的是：


\begin{center}\begin{tabular}{>{\raggedright\arraybackslash}p{\dimexpr(\linewidth-8\tabcolsep)/4\relax}>{\raggedright\arraybackslash}p{\dimexpr(\linewidth-8\tabcolsep)/4\relax}>{\raggedright\arraybackslash}p{\dimexpr(\linewidth-8\tabcolsep)/4\relax}>{\raggedright\arraybackslash}p{\dimexpr(\linewidth-8\tabcolsep)/4\relax}}
\textbf{A.} 社会效益 & \textbf{B.} 经济效益 & \textbf{C.} 繁荣文化市场 & \textbf{D.} 创新文化体制\\[.35em]\end{tabular}\end{center}

7.近年来，为了缩小我国居民在收入分配方面的差距，党和政府做出了巨大的努力，如提高个税起征点、提高企业退休人员基本养老金、提高国家扶贫标准和城乡低保补助水平等，这些举措体现了：


\begin{center}\begin{tabular}{>{\raggedright\arraybackslash}p{\dimexpr(\linewidth-4\tabcolsep)/2\relax}>{\raggedright\arraybackslash}p{\dimexpr(\linewidth-4\tabcolsep)/2\relax}}
\textbf{A.} 初次分配注重效率 & \textbf{B.} 再次分配更加注重公平\\[.35em]
\textbf{C.} 劳动报酬在初次分配中比重提高 & \textbf{D.} 各种生产要素都能按贡献参与分配\\[.35em]\end{tabular}\end{center}

8.党的十八大报告提出，为确保实现全面建成小康社会的宏伟目标，到2020 年，在实现国内生产总值比2010年翻一番的同时，还要实现翻一番的是：


\begin{center}\begin{tabular}{>{\raggedright\arraybackslash}p{\dimexpr(\linewidth-8\tabcolsep)/4\relax}>{\raggedright\arraybackslash}p{\dimexpr(\linewidth-8\tabcolsep)/4\relax}>{\raggedright\arraybackslash}p{\dimexpr(\linewidth-8\tabcolsep)/4\relax}>{\raggedright\arraybackslash}p{\dimexpr(\linewidth-8\tabcolsep)/4\relax}}
\textbf{A.} 城乡居民人均收入 & \textbf{B.} 城乡居民可支配收入 & \textbf{C.} 国民收入 & \textbf{D.} 财政收入\\[.35em]\end{tabular}\end{center}

9.甲午战争后，维新运动迅速兴起。针对洋务派提出的“中体西用”的方针，维新派指出，体与用是不可分的。中学有中学的“体”与“用”，西学有西学的“体”与“用”，把中学之“体”与西学之“用”凑在一起，就如同让“牛体”产生“马用”一样荒谬。维新派与洋务派分歧的实质是：


\begin{center}\begin{tabular}{>{\raggedright\arraybackslash}p{\dimexpr(\linewidth-4\tabcolsep)/2\relax}>{\raggedright\arraybackslash}p{\dimexpr(\linewidth-4\tabcolsep)/2\relax}}
\textbf{A.} 要不要社会革命 & \textbf{B.} 要不要以革命手段推翻清王朝\\[.35em]
\textbf{C.} 要不要在中国兴办近代企业 & \textbf{D.} 要不要学习西方的政治制度与思想文化\\[.35em]\end{tabular}\end{center}

10.1948 年10 月2 日，刘少奇同志在同华北记者团谈话时，讲了一个希腊神话故事：巨人安泰是地神之子，他在同对手搏斗时，只要身不离地，就能从大地母亲那里不断吸取力量，所向无敌；但是，只要他一离开大地，就会毫无力量。他的对手赫拉克勒斯发现了他的这一特征，在一次搏斗中突然把他举到半空中将他扼死。刘少奇借用这一神话故事强调中国共产党始终要：


\begin{center}\begin{tabular}{>{\raggedright\arraybackslash}p{\dimexpr(\linewidth-4\tabcolsep)/2\relax}>{\raggedright\arraybackslash}p{\dimexpr(\linewidth-4\tabcolsep)/2\relax}}
\textbf{A.} 坚持理论联系实际 & \textbf{B.} 保持党的方针政策的正确\\[.35em]
\textbf{C.} 保持对敌人的高度警惕 & \textbf{D.} 保持同人民群众的血肉联系\\[.35em]\end{tabular}\end{center}

11.全面提高公民道德素质，要坚持依法治国和以德治国相结合，加强社会公德、职业道德、家庭美德、个人品德教育，弘扬中华传统美德，弘扬时代新风。下列选项中，既是道德规范，又是法律原则的是：


\begin{center}\begin{tabular}{>{\raggedright\arraybackslash}p{\dimexpr(\linewidth-8\tabcolsep)/4\relax}>{\raggedright\arraybackslash}p{\dimexpr(\linewidth-8\tabcolsep)/4\relax}>{\raggedright\arraybackslash}p{\dimexpr(\linewidth-8\tabcolsep)/4\relax}>{\raggedright\arraybackslash}p{\dimexpr(\linewidth-8\tabcolsep)/4\relax}}
\textbf{A.} 爱岗敬业 & \textbf{B.} 诚实守信 & \textbf{C.} 助人为乐 & \textbf{D.} 勤俭持家\\[.35em]\end{tabular}\end{center}

12.我们要大力弘扬的时代精神是当代中国人民精神风貌的集中体现，是激发社会创造活力的强大力量。时代精神内涵十分丰富，其核心是：


\begin{center}\begin{tabular}{>{\raggedright\arraybackslash}p{\dimexpr(\linewidth-8\tabcolsep)/4\relax}>{\raggedright\arraybackslash}p{\dimexpr(\linewidth-8\tabcolsep)/4\relax}>{\raggedright\arraybackslash}p{\dimexpr(\linewidth-8\tabcolsep)/4\relax}>{\raggedright\arraybackslash}p{\dimexpr(\linewidth-8\tabcolsep)/4\relax}}
\textbf{A.} 国际主义 & \textbf{B.} 集体主义 & \textbf{C.} 改革创新 & \textbf{D.} 开拓进取\\[.35em]\end{tabular}\end{center}

13.个体的人生活动不仅具有满足自我需要的价值属性，还必然包含着满足社会需要的价值属性。个人的需要能不能从社会中得到满足，在多大程度上得到满足，取决于他的：


\begin{center}\begin{tabular}{>{\raggedright\arraybackslash}p{\dimexpr(\linewidth-8\tabcolsep)/4\relax}>{\raggedright\arraybackslash}p{\dimexpr(\linewidth-8\tabcolsep)/4\relax}>{\raggedright\arraybackslash}p{\dimexpr(\linewidth-8\tabcolsep)/4\relax}>{\raggedright\arraybackslash}p{\dimexpr(\linewidth-8\tabcolsep)/4\relax}}
\textbf{A.} 社会影响 & \textbf{B.} 社会价值 & \textbf{C.} 社会地位 & \textbf{D.} 社会理想\\[.35em]\end{tabular}\end{center}

14.“和为贵”是中华民族的传统美德，采用调解的方法解决纠纷，有利于社会和谐。调解可以在诉讼程序外进行，也可以在诉讼程序内进行，诉讼中调解是指：


\begin{center}\begin{tabular}{>{\raggedright\arraybackslash}p{\dimexpr(\linewidth-8\tabcolsep)/4\relax}>{\raggedright\arraybackslash}p{\dimexpr(\linewidth-8\tabcolsep)/4\relax}>{\raggedright\arraybackslash}p{\dimexpr(\linewidth-8\tabcolsep)/4\relax}>{\raggedright\arraybackslash}p{\dimexpr(\linewidth-8\tabcolsep)/4\relax}}
\textbf{A.} 人民调解 & \textbf{B.} 行政调解 & \textbf{C.} 司法调解 & \textbf{D.} 仲裁调解\\[.35em]\end{tabular}\end{center}

15.2012 年6 月27 日，中国宣布在南海地区对外开放九个海上区块，供与外国公司合作勘探开发。此外，海南省宣布将西沙群岛的四个区域划定为文化遗产保护区。这些决定连同设立三沙市，构成了中国加强在南海地区维护主权的“组合行动”。中国在南海问题上的一贯主张是：


\begin{center}\begin{tabular}{>{\raggedright\arraybackslash}p{\dimexpr(\linewidth-4\tabcolsep)/2\relax}>{\raggedright\arraybackslash}p{\dimexpr(\linewidth-4\tabcolsep)/2\relax}}
\textbf{A.} 使南海问题长期化 & \textbf{B.} 推动南海问题国际化\\[.35em]
\textbf{C.} 搁置争议，共同开发 & \textbf{D.} 建立“南海和平自由友谊合作区”\\[.35em]\end{tabular}\end{center}

16.2012 年9 月25 日，第67 届联合国大会一般性辩论在纽约联合国总部开幕。针对错综复杂的国际形势和此起彼伏的地区动荡及热点问题，本次辩论的主题是：


\begin{center}\begin{tabular}{>{\raggedright\arraybackslash}p{\dimexpr(\linewidth-4\tabcolsep)/2\relax}>{\raggedright\arraybackslash}p{\dimexpr(\linewidth-4\tabcolsep)/2\relax}}
\textbf{A.} 改善全球经济治理 & \textbf{B.} 携手推动各国普遍安全与共同发展\\[.35em]
\textbf{C.} 增强联合国维和作用 & \textbf{D.} 以和平方式调解或解决国际争端或局势\\[.35em]\end{tabular}\end{center}

\subsection*{二、多项选择题}

17.马克思主义是关于无产阶级和人类解放的科学，实现共产主义是全人类解放的根本体现。人类解放包括：


\begin{center}\begin{tabular}{>{\raggedright\arraybackslash}p{\dimexpr(\linewidth-4\tabcolsep)/2\relax}>{\raggedright\arraybackslash}p{\dimexpr(\linewidth-4\tabcolsep)/2\relax}}
\textbf{A.} 从自然界的压迫下解放出来 & \textbf{B.} 从客观规律的制约下解放出来\\[.35em]
\textbf{C.} 从旧的社会关系的束缚下解放出来 & \textbf{D.} 从旧的传统观念的禁锢下解放出来\\[.35em]\end{tabular}\end{center}

18.唯物史观第一次科学地解决了历史创造者的问题，认为人民群众是历史的创造者。人民群众：


\begin{center}\begin{tabular}{>{\raggedright\arraybackslash}p{\dimexpr(\linewidth-4\tabcolsep)/2\relax}>{\raggedright\arraybackslash}p{\dimexpr(\linewidth-4\tabcolsep)/2\relax}}
\textbf{A.} 从量上说是指社会人口的绝大多数 & \textbf{B.} 从质上说是对社会历史发展起推动作用的人们\\[.35em]
\textbf{C.} 在任何历史时期都不包括剥削阶级 & \textbf{D.} 最稳定的主体部分始终是从事物质资料生产的劳动群众及其知识分子\\[.35em]\end{tabular}\end{center}

19.美国导演迈克尔·穆尔的最新纪录片《资本主义：一个爱情故事》问世以来，一直颇受关注。“资本主义”为何与“爱情故事”联系起来呢？穆尔解释说，这是一种“贪欲之爱”。“喜爱财富的人不仅爱他们自己的钱，也爱你口袋中的钱……很多人不敢说出它的名字，真见鬼，就说出来吧。这就是资本主义。”对金钱的“贪欲”之所以与资本主义连为一体，是因为：


\begin{center}\begin{tabular}{>{\raggedright\arraybackslash}p{\dimexpr(\linewidth-4\tabcolsep)/2\relax}>{\raggedright\arraybackslash}p{\dimexpr(\linewidth-4\tabcolsep)/2\relax}}
\textbf{A.} 资本家是人格化的资本 & \textbf{B.} 赚钱体现了人的天然本性\\[.35em]
\textbf{C.} 资本的生命在于不断运动和不断增殖 & \textbf{D.} 追逐剩余价值是资本主义生产方式的绝对规律\\[.35em]\end{tabular}\end{center}

20.伴随着生产力发展、科技进步及阶级关系调整，当代资本主义社会的劳资关系和分配关系发生了很大变化。其中资本家及其代理人为缓和劳资关系所采取的激励制度有：


\begin{center}\begin{tabular}{>{\raggedright\arraybackslash}p{\dimexpr(\linewidth-8\tabcolsep)/4\relax}>{\raggedright\arraybackslash}p{\dimexpr(\linewidth-8\tabcolsep)/4\relax}>{\raggedright\arraybackslash}p{\dimexpr(\linewidth-8\tabcolsep)/4\relax}>{\raggedright\arraybackslash}p{\dimexpr(\linewidth-8\tabcolsep)/4\relax}}
\textbf{A.} 职工参与决策制度 & \textbf{B.} 职工终身雇佣制度 & \textbf{C.} 职工选举管理者制度 & \textbf{D.} 职工持股制度\\[.35em]\end{tabular}\end{center}

21.自第一个社会主义国家建立以来，社会主义事业的发展并不是一帆风顺的。社会主义发展道路的多样性以及发展过程中的前进性和曲折性的实践告诉我们：


\begin{center}\begin{tabular}{>{\raggedright\arraybackslash}p{\dimexpr(\linewidth-4\tabcolsep)/2\relax}>{\raggedright\arraybackslash}p{\dimexpr(\linewidth-4\tabcolsep)/2\relax}}
\textbf{A.} 坚持社会主义，不等于要坚持某种单一的社会主义模式 & \textbf{B.} 发展社会主义，不等于不学习西方资本主义的文明成果\\[.35em]
\textbf{C.} 改革或抛弃某种社会主义模式，不等于改掉或抛弃社会主义 & \textbf{D.} 某种社会主义模式的失败，不等于整个社会主义事业的成败\\[.35em]\end{tabular}\end{center}

22.党的十八大把科学发展观同马克思列宁主义、毛泽东思想、邓小平理论、“三个代表”重要思想一道确立为党必须长期坚持的指导思想。科学发展观是：


\begin{center}\begin{tabular}{>{\raggedright\arraybackslash}p{\dimexpr(\linewidth-4\tabcolsep)/2\relax}>{\raggedright\arraybackslash}p{\dimexpr(\linewidth-4\tabcolsep)/2\relax}}
\textbf{A.} 中国革命、建设、改革经验的科学总结 & \textbf{B.} 中国特色社会主义理论体系的最新成果\\[.35em]
\textbf{C.} 中国共产党集体智慧的结晶 & \textbf{D.} 指导党和国家全部工作的强大思想武器\\[.35em]\end{tabular}\end{center}

23.国务院新闻办公室发表的《西藏和平解放60 年》白皮书指出:目前，西藏共有各类宗教活动场所1700 余处，僧尼约4.6 万人，藏传佛教特有的活佛转世的传承方式得到充分尊重，寺庙学经、辩经、受戒、灌顶、修行等传统宗教活动和寺庙学经考核晋升学位活动正常进行，每年到拉萨朝佛敬香的信教群众达百万人次以上。上述事实表明，我国的宗教政策得到了充分贯彻。我国宗教政策主要有：


\begin{center}\begin{tabular}{>{\raggedright\arraybackslash}p{\dimexpr(\linewidth-4\tabcolsep)/2\relax}>{\raggedright\arraybackslash}p{\dimexpr(\linewidth-4\tabcolsep)/2\relax}}
\textbf{A.} 积极引导宗教与社会主义社会相适应 & \textbf{B.} 尊重和保护公民的宗教信仰自由\\[.35em]
\textbf{C.} 独立自主自办教会 & \textbf{D.} 依法管理宗教事务\\[.35em]\end{tabular}\end{center}

24.统筹区域发展，促进区域协调发展，是我国经济社会发展的一个重点原则。坚持这一原则有利于：


\begin{center}\begin{tabular}{>{\raggedright\arraybackslash}p{\dimexpr(\linewidth-4\tabcolsep)/2\relax}>{\raggedright\arraybackslash}p{\dimexpr(\linewidth-4\tabcolsep)/2\relax}}
\textbf{A.} 扩大内需，拉动经济增长 & \textbf{B.} 区域间优势互补，促进经济共同发展\\[.35em]
\textbf{C.} 不同区域人民共享改革发展成果 & \textbf{D.} 生产要素在区域间合理流动和配置\\[.35em]\end{tabular}\end{center}

25.1992 年初，邓小平在南方谈话中指出：“社会主义的本质，是解放生产力，发展生产力，消除剥削，消除两极分化，最终达到共同富裕。”这一概括对社会主义传统认识的突破主要体现在：


\begin{center}\begin{tabular}{>{\raggedright\arraybackslash}p{\dimexpr(\linewidth-4\tabcolsep)/2\relax}>{\raggedright\arraybackslash}p{\dimexpr(\linewidth-4\tabcolsep)/2\relax}}
\textbf{A.} 破除了脱离生产力水平抽象谈论社会主义的认识 & \textbf{B.} 否定了社会主义必须坚持公有制和按劳分配原则的认识\\[.35em]
\textbf{C.} 摆脱了长期以来忽视建设社会主义根本目的和目标的认识 & \textbf{D.} 纠正了把社会主义本质等同于社会主义具体做法的认识\\[.35em]\end{tabular}\end{center}

26.PM2.5（细粒颗粒），这个大家原本很陌生的专有名词，因为2011 年10 月我国多地灰霾天气造成严重大气污染而迅速成为社会热词。2012 年2 月修订的《环境空气质量标准》增加了PM2.5 指标，该指标随后又被写入政府工作报告。这既折射出当前我国环境污染的严重性，同时也反映了党和政府治理环境污染、建设生态文明的决心。党的十八大进一步把生态文明建设融入中国特色社会主义建设五位一体的总布局。十八大提出的生态文明建设新要求是：


\begin{center}\begin{tabular}{>{\raggedright\arraybackslash}p{\dimexpr(\linewidth-4\tabcolsep)/2\relax}>{\raggedright\arraybackslash}p{\dimexpr(\linewidth-4\tabcolsep)/2\relax}}
\textbf{A.} 加大自然生态系统和环境保护力度 & \textbf{B.} 全面促进资源节约\\[.35em]
\textbf{C.} 优化国土空间开发格局 & \textbf{D.} 加强生态文明制度建设\\[.35em]\end{tabular}\end{center}

27.1925 至1927 年的大革命，规模宏伟，内涵丰富。与辛亥革命相比较，其不同点在于：


\begin{center}\begin{tabular}{>{\raggedright\arraybackslash}p{\dimexpr(\linewidth-4\tabcolsep)/2\relax}>{\raggedright\arraybackslash}p{\dimexpr(\linewidth-4\tabcolsep)/2\relax}}
\textbf{A.} 它广泛而深刻地发动了工农群众 & \textbf{B.} 它的主要斗争形式是武装斗争\\[.35em]
\textbf{C.} 它的革命对象是帝国主义和封建军阀 & \textbf{D.} 它是在以国共合作为基础的统一战线的组织形式下进行的\\[.35em]\end{tabular}\end{center}

28.1931 年1 月至1935 年1 月，以王明为代表的“左”倾错误给中国革命带来严重危害，其主要错误有：


\begin{center}\begin{tabular}{>{\raggedright\arraybackslash}p{\dimexpr(\linewidth-4\tabcolsep)/2\relax}>{\raggedright\arraybackslash}p{\dimexpr(\linewidth-4\tabcolsep)/2\relax}}
\textbf{A.} 排斥和打击中间势力 & \textbf{B.} 将反帝反封建与反资产阶级并列\\[.35em]
\textbf{C.} 集中力量攻打大城市 & \textbf{D.} 主张“一切经过统一战线”\\[.35em]\end{tabular}\end{center}

29.抗日战争是近代以来中华民族反抗外敌入侵第一次取得完全胜利的民族解放战争。中国赢得抗日战争胜利的主要原因是：


\begin{center}\begin{tabular}{>{\raggedright\arraybackslash}p{\dimexpr(\linewidth-4\tabcolsep)/2\relax}>{\raggedright\arraybackslash}p{\dimexpr(\linewidth-4\tabcolsep)/2\relax}}
\textbf{A.} 中国共产党发挥了中流砥柱的作用 & \textbf{B.} 中国的国力空前强大\\[.35em]
\textbf{C.} 得到了国际反法西斯力量的同情与支持 & \textbf{D.} 中国实现了空前的民族觉醒和民族团结\\[.35em]\end{tabular}\end{center}

30.以毛泽东同志为核心的党的第一代中央领导集体带领全党全国各族人民完成了新民主主义革命，进行了社会主义改造，确立了社会主义基本制度。这一基本制度的确立：


\begin{center}\begin{tabular}{>{\raggedright\arraybackslash}p{\dimexpr(\linewidth-4\tabcolsep)/2\relax}>{\raggedright\arraybackslash}p{\dimexpr(\linewidth-4\tabcolsep)/2\relax}}
\textbf{A.} 为当代中国一切发展进步奠定了根本政治前提和制度基础 & \textbf{B.} 是中国历史上最深刻最伟大的社会变革\\[.35em]
\textbf{C.} 标志着马克思主义同中国实际“第二次结合”的完成 & \textbf{D.} 使广大劳动人民真正成为国家的主人\\[.35em]\end{tabular}\end{center}

31.一位社会学家发现大楼的一块玻璃坏了，起初他没太当回事，没过多久，他发现许多处窗户都破损了。经过调研后，他得出结论：一样东西如果有一点破损，人们就会有意无意地加快它的破损速度；一样东西如果完好无损，或是及时维护，人们就会精心地护理。这就是著名的“破窗定律”。下列关于道德修养的名言与“破窗定律”内涵相近的是：


\begin{center}\begin{tabular}{>{\raggedright\arraybackslash}p{\dimexpr(\linewidth-4\tabcolsep)/2\relax}>{\raggedright\arraybackslash}p{\dimexpr(\linewidth-4\tabcolsep)/2\relax}}
\textbf{A.} 非知之难，行之惟难；非行之难，终之斯难 & \textbf{B.} 善不可谓小而无益，不善不可谓小而无伤\\[.35em]
\textbf{C.} 小善虽无大益，而不可不为；细恶虽无近祸，而不可不去也 & \textbf{D.} 见贤思齐焉，见不贤而内自省也\\[.35em]\end{tabular}\end{center}

32.“可上九天揽月，可下五洋捉鳖”。2012 年6 月24 日，在航天员的手动操控下，我国“神舟九号”宇宙飞船与“天宫一号”交会对接成功，向未来建设宇宙空间站迈出坚实一步。就在同一天，我国的“蛟龙号”载人潜水器成功下潜7020 米的深度，再次创造历史记录。这两项成就的里程碑意义在于：


\begin{center}\begin{tabular}{>{\raggedright\arraybackslash}p{\dimexpr(\linewidth-4\tabcolsep)/2\relax}>{\raggedright\arraybackslash}p{\dimexpr(\linewidth-4\tabcolsep)/2\relax}}
\textbf{A.} 完成了“近地空间长期载人飞行”和“南海深部计划” & \textbf{B.} 丰富了人类认识和开发宇宙的梦想和实践\\[.35em]
\textbf{C.} 为中国在太空和海底参与国际合体提供了更多机会 & \textbf{D.} 为中国在太空和大洋这两个尚未充分利用空间的活动创造了新的可能性\\[.35em]\end{tabular}\end{center}

33.2012 年6 月6 日至7 日，上海合作组织成员国元首理事会第十二次会议在北京举行。这是该组织发展进入第二个十年的首次峰会，与会各国元首就深化成员国友好合作以及重大国际和地区问题深入交换了意见，达成了新的重要共识。此次峰会的主要成果有：


\begin{center}\begin{tabular}{>{\raggedright\arraybackslash}p{\dimexpr(\linewidth-4\tabcolsep)/2\relax}>{\raggedright\arraybackslash}p{\dimexpr(\linewidth-4\tabcolsep)/2\relax}}
\textbf{A.} 首次制定了中长期战略规划 & \textbf{B.} 签订了军事同盟条约\\[.35em]
\textbf{C.} 签署了首个人文合作宣言 & \textbf{D.} 增加了新的观察员国和对话伙伴国\\[.35em]\end{tabular}\end{center}

\subsection*{三、分析题}

34.结合材料回答问题：


\subsection*{材料1}

小学老师雷夫·爱斯奎斯在其所著的热门教育畅销书《第56 号教室的奇迹》中讲过这样一个故事：一位从事特殊教育的优秀教师获得一个宝贵的签名球，上面有美国著名棒球队---红袜队1967 年全体队员的签名。这些球员都是他的偶像，对这样一个签名球，这位教师别提有多珍爱了。当年幼的儿子找他一起玩球时，他理所当然地警告儿子：绝对不能拿签名球来玩。儿子问他理由，他觉得儿子还太小，对球队和球员一无所知，说多了儿子也不会明白。于是，他没有解释原委，只对儿子说，不能用那颗球，是因为“球上写满了字”。过了几天，儿子又找他一起玩球。当老爸再次提醒儿子不要拿写满字的球来玩时，小男孩满不在乎说：我已经把问题解决了。爸爸问怎么回事，儿子说：我把球上所有的字都擦掉了。老爸气的想痛打儿子，但他转念一想，觉得儿子根本没有做错事，因为自己并没有告诉儿子上面的字有什么意义。从那天起，他无论去什么地方都带着那颗空白的签名球。这颗球提醒他，不管是教导学生还是子女，一定要时时从孩子的角度去看事情。不论家长还是教师，常常用成人的眼光看待孩子，用成人的思维理解孩子，用成人的标准要求孩子。岂不知，从孩子的角度看事情，用孩子的眼睛看世界，正是儿童教育应当遵循的基本规律。摘编自《人民日报》（2012 年3 月16 日）


\subsection*{材料2}

某大学一研究生凭借着设计“醒目药瓶”，摘得了素有“设计界奥斯卡”美誉的2011 年度“国际红点奖”概念设计类奖。在她提供的设计图上，常见的塑料药瓶盖的顶上一圈，变身为一块圆圆的玻璃。“这是一面凹凸镜，有放大的功能”。他解释说，有了这个药瓶盖，老年人不需要带上老花镜来区别药的类别、服用量等。她的灵感来源于生活中对中老年人群体的关注。有一天，有位老人要吃药，可是药瓶上的字太小了，原本挂在脖子上的老花镜却不见了，急得这位老人团团转。就这样，该同学很长一段时间沉浸在老人世界中，突然有一天灵感迸发，想到“醒目药瓶”这个点子。有了灵感后，从设计，到写英文翻译说明，再到制作动画，一共才三天时间。也许有人要问，这样的设计看上去很简单，为什么能拿“国际红点奖”呢？他坦言，设计很简单，关键在于设计前把自己想像成老人。这一设计胜在实用，按照测算，不会给药瓶本身带来额外的成本，推广起来很容易，使用方便。“希望将来这款设计能推向市场，让更多人得到帮助。”这位研究生说她没有想当名人的“野心”。只期望能从生活中的小处入手，用自己的设计改变生活，让生活更加美好。正如“红点”主席Peter Zec 博士在颁奖晚会上说的那样：从同学们优秀的设计中，他高兴地看到的是他们所描绘的未来更加美好的世界。摘编自《扬子晚报》（2012 年3 月17 日）


(1)分析“用孩子的眼睛看世界”和“设计前把自己想象成老人”两事例所体现的认识主体的能动作用。（6 分）


(2)“用自己的设计改变生活，让生活更加美好”对我们从事实践活动有何意义？（4 分）


35.结合材料回答问题：


\subsection*{材料1}

浙江武义县后陈村是全县经济条件较好的村，但由于财务使用不透明，村民从上世纪90 年代末以来连续向县纪委、街道反映村里的问题，却长期得不到彻底解决。2002 年和2003 年，连续两任村支书因经济问题被查处。2004 年初，村里有1000 亩土地被征用，获补偿金1900 多万元。在人均发放7000 元后，还剩1000 多万元。如何处理这些集体资产成为村内一个难题。为帮助村民寻找从根本上解决问题的办法，由县纪委牵头组建的村务监督改革指导小组进驻后陈村。指导组在大量听取村民意见的基础上，决定组建一个相对独立于村委会及村党支部的监督委员会，真正能从根本上让村民有效制约村干部的权力。6 月，后陈村在海选村委会基础上由群众选举产生了全国第一个村级民主监督组织——村务监督委员会，与村党支部、村委会一起称为“三委会”。监委会成立不久，即对村里两口池塘的承包进行了全程监督，结果每口池塘三年的承包价从2.8 万升至


5.8 万元。2005 年，后陈村举行垃圾清运投标会。村委会主任主持会议，监委会主任到场监督，不到5 分钟结果就出来了。没有人对投标公正性表示质疑。监委会成立后，后陈村每年的创收情况，包括出租土地给广告公司做广告牌、旧粮站出租、经营沙场、村留土地上的杉树出售，以及向上级部门申请到的资金补助等，每一笔都要经过监委会审核后公布。2004年当年的招待开支是23909 元，比前些年下降近一半，村干部再不能拿着发票随便报销了。在村财务公开栏前，村民告诉记者：“过去简直是胡来，集体的钱像是干部自己的，现在不一样了。”监委会主任说：“我的职责就是看他们有没有按程序办事，有没有搞暗箱操作。”监委会成立以来，后陈村的固定收入逐年增加，村干部连续8 年实现零违纪，村民连续8 年实现零上访。村两委已顺利完成了3 次换届。最近的一次换届，村两委成员一个没动，全部高票当选，一次通过。目前，浙江省3 万多个行政村，村村建立了村务监督委员会，实现了村级监督组织全覆盖。村务监督委员会这一有效而不需要太大监督成本的权力制衡制度，对建立乡村“阳光权力体系”，共建和谐社会带来重要启迪。村务监督委员会这一制度创新已被写进《中华人民共和国村民委员会组织法》，并在全国推行。摘编自《人民日报》（2012 年5 月14 日）等


(1)后陈村是如何通过制度创新来保障权力在阳光下运行的？（5 分）


(2)全国第一个村务监督委员会的建立对推进基层民主制度建设有何启示？（5 分）


36.结合材料回答问题：


\subsection*{材料1}

1910 年，上海人陆士谔在幻想小说《新中国》里记载了一个神奇的梦。梦中主人公随时光穿梭，看到“万国博览会”在上海浦东举行，为方便市民参观，上海滩建成了浦东大铁桥和越江隧道，还造了地铁，工厂中的机器有鬼斧神工之妙，租界的治外法权已经收回，汉语成了世界通用的流行语言……最后梦中人一跤跌醒，却言道：“休说是梦，到那时，真有这景象也未可知。”1920 年，孙中山先生完成《建国方略》一书，书中提出了修建三峡水利、建设高原铁路系统等宏伟设想，构想了工厂遍地、机器轰鸣、高楼大厦矗立城乡、火车轮船繁忙往返的现代化景象，描绘了“万众一心，急起直追，以我五千年文明优秀之民族，应世界之潮流，而建设一政治最修明、人民最安乐之国家”的愿景。1935 年，方志敏在《可爱的中国》中写道：“中国一定有个可赞美的光明前途……到那时，到处都是活跃的创造，到处都是日新月异的进步，欢歌将代替了悲叹，笑脸将代替了哭脸，富裕将代替了贫穷，康健将代替了疾苦，智慧将代替了愚昧，友爱将代替了仇杀，生之快乐将代替了死之悲哀，明媚的花园，将代替了凄凉的荒地！这时，我们民族就可以无愧色的立在人类的面前，而生育我们的母亲，也会最美丽地装饰起来，与世界上各位母亲平等的携手了。” “这么光荣的一天，决不在辽远的将来，而在很近的将来。”摘编自《经济日报》（2012 年12 月2 日）、《方志敏文集》


\subsection*{材料2}

2012 年11 月29 日，中共中央总书记习近平到国家博物馆参观《复兴之路》展览。在十九世纪末列强割占领土、设立租借地、划定势力范围示意图前，在鸦片战争期间虎门的大炮前，在反映辛亥革命的文物和照片前，在《共产党宣言》第一个中文全译本前，在《中国共产党的第一个纲领》等反映中国共产党成立的文物和照片前，在李大钊狱中亲笔自述前，在中华人民共和国第一面五星红旗前，在党的十一届三中全会照片前，习近平不时停下脚步，认真观看，仔细询问和了解有关情况。在参观过程中，习近平发表了重要讲话。他指出，每个人都有理想和追求，都有自己的梦想。实现中华民族伟大复兴，就是中华民族近代以来最伟大的梦想。中华民族的昨天，可以说是“雄关漫道真如铁”；中华民族的今天，正可谓“人间正道是沧桑”；中华民族的明天，可以说是“长风破浪会有时”。经过鸦片战争以来170 多年的持续奋斗，中华民族伟大复兴展现出光明的前景。现在，我们比历史上任何时候都更接近中华民族伟大复兴的目标，比历史上任何时期都更有信心、有能力实现这个目标。摘编自《人民日报》（2012 年11 月30 日）


(1)为什么“实现中华民族伟大复兴，就是中华民族近代以来最伟大的梦想”？（4 分）


(2)为什么说“现在，我们比历史上任何时期都更接近中华民族伟大复兴的目标”？（6 分）


37.结合材料回答问题：某图书馆向所有读者免费开放。乞丐、拾荒者、衣衫破旧的民工小心翼翼进来了，无人阻拦。于是他们便堂而皇之地在馆内读书、看报。有读者对此表示不满，向馆长抱怨说：“图书馆是大雅之堂，如果允许乞丐和拾荒者进入阅读，就是对其他读者的不尊重。”馆长回答说：“我无权拒绝他们入内阅读，但你有权选择离开。”此事被发在微博上，顿时触动了社会的神经，引发人们对人文精神的关注和思考。中央电视台等主流媒体对此进行了报道，一场公共图书馆办馆理念的大讨论由此引发。公共图书馆一向更愿意向体面的文化人敞开，常在门口赫然告示：“衣冠不整，谢绝入内！”把读者分为三六九等，拒绝部分人入内，其公益性大打折扣。而该馆馆长希望图书馆“成为每一个读者的天堂”，“无论任何人，只要进了图书馆，在知识面前都享有同等权利，不得有高低贵贱之分。”为此，该馆在全国同行中率先推出免证阅读制度，任何人进馆借阅书籍都不需要证件和费用，以体现人道、人文的公共图书馆办馆理念和人性化的服务。对于图书馆实行免费开放可能带来的问题，该馆有关负责人感触颇深：“自图书馆实行零门槛后，我们不仅没有感到压力增加，反而感觉开放的时间越长，不尊重这种权利的读者越少，我们和读者都被这种和谐的环境所改变。至于进管要先洗手，馆内并没有硬性规定，耳濡目染的时间长了，谁也会自觉地洗手，然后再阅读。”“如果有天堂，天堂应该是图书馆的模样。”这是文学大师、曾担任阿根廷国立图书馆馆长的博尔赫斯的一句名言。该图书馆向乞丐和拾荒者免费开放，不啻一轮明亮的太阳，让乞丐和拾荒者在得到温暖的同时，也净化我们的心灵。摘编自《中国青年》（2011 年第5 期）、《光明日报》（2012 年5 月10 日）


(1)从法律角度如何理解“我无权拒绝他们入内阅读，但你有权选择离开”？（5 分）


(2)图书馆向乞丐和拾荒者免费开放对我们处理人际关系有何启示？（5 分）


38.阅读下列材料：


\subsection*{材料1}

2010 年5 月，美国发布奥巴马上台后第一个《四年防务评估报告》和《国家安全战略报告》。报告指出，美国国家利益由安全、繁荣、价值和国际秩序四个方面组成；美国通过对这些利益的追求，实现国家复兴和全球领导地位；相比世界其他地区，亚洲是美国最有所作为的地区。同时，报告将东南亚国家归为三个类别：“正式盟友”、“战略伙伴”和“可预期的战略伙伴”，虽然，美国准备让东盟所有国家成为美国的盟友或是伙伴。美国卡内基国际和平基金会副总裁包道格说：过去10 年来，中国在东南亚地区不断扩大其利益，取得了有效的成果，这是美国没有做到的。摘编自新华网（2011 年5 月28 日）


\subsection*{材料2}

2011 年10 月以来，美国高层不断访问亚太地区的国家，参加相关的国际会议。10 月下旬，国防部长帕内塔出访印尼，日本和韩国，强调美国将加大在亚太的军事部署。11 月下旬到12 月初，国务卿希拉里·克林顿，先后访问菲律宾、泰国，并对缅甸进行了“历史性访问”，这是1955 年以来美国国务卿首次访问缅甸。同月，奥巴马也展开亚太之行。他参加在夏威夷举行的亚太经合组织第十九次领导人非正式会议，随后出访澳大利亚并前往印尼出席东亚多边峰会，成为参加东亚峰会的首位美国总统。奥巴马在亚太之行中高调宣示，美国是“太平洋大国”，将“留驻”亚太，通过“坚持核心原则”和与盟友及伙伴的紧密合作，在“塑造”亚太地区未来中发挥更大、更长远的作用。2012 年11 月8 日，奥巴马连任成功不到48 小时，即宣布他的首次出访选在东南亚。17 日至20 日，他不仅访问了泰国，而且对缅甸和柬埔寨进行了历史性首访。在这次东南亚之行中，奥巴马再次强调，美国是一个太平洋国家，亚太地区对美国创造就业机会以及塑造其安全与繁荣至关重要。摘编自《人民日报》（2011 年12 月23 日）、《参考消息》（2012 年11 月21 日）等


\subsection*{材料3}

第十一届香格里拉对话暨亚洲安全会议于2012 年6 月1 日至3 日在新加坡举行。会议期间，美国国防部长帕内塔发表了题为《美国对亚太的再平衡》的演讲，重点阐释了“亚太再平衡战略”的军事计划，其中包括2020 年前在亚太地区保持6 个航母舰队，以及将60\%的海军力量部署到亚太地区。国际舆论普遍怀疑，美国战略重心向亚太地区转移的意图是为了遏制中国。摘编自新华网（2012 年6 月4 日）


（1）美国将其全球战略重心转向亚太的原因何在？（5 分）


（2）如何看待美国战略重心东移对中国周边安全的影响？（5 分）

\end{document}
